\chapter{Recuit simulé et recuit quantique}\label{ch:recuit}
\iterbanner{V0}

Deux façons de chercher l'état fondamental en « refroidissant » progressivement : avec du hasard thermique (recuit \emph{simulé}, classique) ou avec des fluctuations quantiques (recuit \emph{quantique}).

\section{Le paysage d'énergie}
On imagine les $2^n$ chaînes comme les points d'un paysage, où deux chaînes sont voisines si elles diffèrent d'un bit, et l'altitude est l'énergie. Une descente « gloutonne », qui retourne un bit dès que l'énergie baisse, s'arrête dans le premier \textbf{minimum local} venu.

\begin{figure}[htbp]
\centering
\begin{tikzpicture}[x=1.3cm,y=1.1cm]
\draw[->,cmuted] (0,-0.2) -- (0,3.1) node[above,cink,font=\small] {énergie};
\draw[->,cmuted] (0,-0.2) -- (9.3,-0.2) node[below left,cink,font=\small] {configurations (chaînes de bits)};
\draw[cink,very thick] plot[smooth,tension=0.7] coordinates {(0.2,2.3) (1.2,1.6) (2.2,1.2) (3.2,1.7) (4.3,2.7) (5.3,1.8) (6.4,0.25) (7.4,1.0) (8.4,2.4) (9.0,2.6)};
\fill[cblue] (2.2,1.2) circle (2.2pt);
\node[font=\small,anchor=north east] at (2.1,1.1) {minimum local};
\node[font=\small,anchor=west] at (6.6,0.2) {minimum global};
\draw[-{Stealth},corange,thick] (2.3,1.35) .. controls (3.4,3.6) and (4.9,3.6) .. (5.7,1.6);
\node[corange,font=\small,above] at (4.0,3.15) {saut thermique (recuit simulé)};
\draw[-{Stealth},cblue,thick,dashed] (2.4,1.25) -- (5.75,1.25);
\node[cblue,font=\small,align=center] at (4.0,0.75) {effet tunnel\\ (recuit quantique)};
\end{tikzpicture}
\caption{Deux façons de quitter un minimum local : passer \emph{par-dessus} la barrière grâce à l'agitation thermique, ou \emph{à travers} grâce aux fluctuations quantiques.}\label{fig:paysage}
\end{figure}

\section{Recuit simulé (classique)}
Inspiré de la métallurgie : un métal refroidi lentement cristallise dans un état de basse énergie. On introduit une \textbf{température} $T$ et on fait une marche aléatoire (algorithme de Metropolis) :
\begin{enumerate}
\item proposer de retourner un bit, ce qui change l'énergie de $\Delta E$ ;
\item accepter si $\Delta E\le0$, sinon accepter avec probabilité $e^{-\Delta E/T}$ ;
\item diminuer lentement $T$ (programme de température, par exemple géométrique).
\end{enumerate}
À $T$ fixée, la marche échantillonne la distribution de Boltzmann $\propto e^{-E(x)/T}$, qui se concentre sur le minimum global quand $T\to0$. Si $T$ baisse trop vite, on reste piégé ; la théorie garantit la convergence pour un refroidissement logarithmique, beaucoup trop lent en pratique.

Retourner le bit $i$ coûte $\Delta E_i=(1-2b_i)\big(Q_{ii}+\sum_{j\neq i}Q_{\min,\max}b_j\big)$. C'est un calcul en $O(n)$, ce qui rend le recuit simulé très rapide (\texttt{quantique/recuit.py}).

\begin{figure}[htbp]
\centering
% figure générée par tools/make_figures.py : recuit_succes
\begin{tikzpicture}
\begin{semilogxaxis}[width=0.7\linewidth,height=5.2cm,axis lines=left,ymin=0,ymax=1.02,
  xlabel={nombre de balayages (pas de température)},ylabel={fraction des 200 lectures},
  legend style={draw=none,font=\small,at={(0.98,0.03)},anchor=south east}]
\addplot[cblue,thick,mark=*] coordinates {(2,0.16) (5,0.33) (10,0.37) (20,0.51) (50,0.445) (100,0.49)};\addlegendentry{minimum exact du QUBO}
\addplot[corange,thick,mark=square*] coordinates {(2,0.435) (5,0.79) (10,0.89) (20,0.94) (50,0.965) (100,0.98)};\addlegendentry{à moins de 0{,}1 count du minimum}
\end{semilogxaxis}
\end{tikzpicture}

\caption{Recuit simulé sur notre QUBO aéro à 9 bits, 200 lectures indépendantes par point. La réussite « exacte » plafonne vers 50~\% : le fondamental et le premier niveau excité sont deux profils physiquement identiques ($m=0$, $p=0{,}25$ ou $0{,}30$), séparés de $0{,}06$ count. Le recuit tombe au hasard dans l'un des deux. À $0{,}1$ count près, la réussite dépasse 95~\% dès 50 balayages.}\label{fig:recuit}
\end{figure}

\section{Le calcul quantique adiabatique}
\paragraph{Idée.} On prépare un système dans l'état fondamental, \emph{facile}, d'un hamiltonien simple $H_B$, puis on le transforme \emph{lentement} en notre $H_C$ :
\[H(s)=(1-s)\,H_B+s\,H_C,\qquad s=t/T\in[0,1],\qquad H_B=-\sum_kX_k .\]
L'état fondamental de $H_B$ est $\ket+^{\otimes n}$, la superposition uniforme de toutes les chaînes (exercice~\ref{ch:hamiltoniens}.2).

\paragraph{Théorème adiabatique.} Si l'évolution est assez lente, le système reste dans l'état fondamental instantané de $H(s)$, et arrive donc dans celui de $H_C$ : la solution. « Assez lent » dépend du plus petit \textbf{écart spectral} (\emph{gap}) $g_{\min}$ entre les deux niveaux les plus bas le long du chemin : il faut typiquement $T\gg1/g_{\min}^2$.

\begin{figure}[htbp]
\centering
% figure générée par tools/make_figures.py : spectre_adiabatique
\begin{tikzpicture}
\begin{axis}[width=0.8\linewidth,height=6.2cm,axis lines=left,xmin=0,xmax=1,ymin=-3.3,ymax=1.3,
  xlabel={$s=t/T$},ylabel={niveaux d'énergie de $H(s)$},
  legend style={draw=none,font=\small,at={(0.02,0.02)},anchor=south west,fill=white,fill opacity=0.85,text opacity=1}]
\addplot[cblue,very thick] coordinates {(0,-3) (0.02,-2.95) (0.04,-2.9) (0.06,-2.851) (0.08,-2.802) (0.1,-2.753) (0.12,-2.704) (0.14,-2.656) (0.16,-2.608) (0.18,-2.56) (0.2,-2.513) (0.22,-2.466) (0.24,-2.419) (0.26,-2.373) (0.28,-2.327) (0.3,-2.282) (0.32,-2.238) (0.34,-2.194) (0.36,-2.152) (0.38,-2.109) (0.4,-2.068) (0.42,-2.028) (0.44,-1.989) (0.46,-1.952) (0.48,-1.916) (0.5,-1.881) (0.52,-1.849) (0.54,-1.818) (0.56,-1.79) (0.58,-1.765) (0.6,-1.742) (0.62,-1.723) (0.64,-1.707) (0.66,-1.695) (0.68,-1.686) (0.7,-1.682) (0.72,-1.681) (0.74,-1.684) (0.76,-1.69) (0.78,-1.7) (0.8,-1.714) (0.82,-1.73) (0.84,-1.75) (0.86,-1.773) (0.88,-1.798) (0.9,-1.826) (0.92,-1.856) (0.94,-1.889) (0.96,-1.924) (0.98,-1.961) (1,-2)};
\addplot[corange,thick] coordinates {(0,-1) (0.02,-1.004) (0.04,-1.008) (0.06,-1.013) (0.08,-1.017) (0.1,-1.021) (0.12,-1.026) (0.14,-1.03) (0.16,-1.035) (0.18,-1.039) (0.2,-1.044) (0.22,-1.048) (0.24,-1.053) (0.26,-1.057) (0.28,-1.062) (0.3,-1.066) (0.32,-1.071) (0.34,-1.075) (0.36,-1.079) (0.38,-1.083) (0.4,-1.087) (0.42,-1.091) (0.44,-1.094) (0.46,-1.097) (0.48,-1.099) (0.5,-1.101) (0.52,-1.101) (0.54,-1.101) (0.56,-1.1) (0.58,-1.098) (0.6,-1.095) (0.62,-1.089) (0.64,-1.083) (0.66,-1.075) (0.68,-1.065) (0.7,-1.054) (0.72,-1.042) (0.74,-1.029) (0.76,-1.015) (0.78,-1.002) (0.8,-0.9892) (0.82,-0.9774) (0.84,-0.9673) (0.86,-0.9593) (0.88,-0.954) (0.9,-0.9518) (0.92,-0.9533) (0.94,-0.9587) (0.96,-0.9682) (0.98,-0.982) (1,-1)};
\addplot[cmuted,thin] coordinates {(0,-1) (0.02,-0.9901) (0.04,-0.9802) (0.06,-0.9705) (0.08,-0.9609) (0.1,-0.9514) (0.12,-0.942) (0.14,-0.9328) (0.16,-0.9238) (0.18,-0.9149) (0.2,-0.9062) (0.22,-0.8977) (0.24,-0.8894) (0.26,-0.8813) (0.28,-0.8735) (0.3,-0.8659) (0.32,-0.8586) (0.34,-0.8515) (0.36,-0.8448) (0.38,-0.8385) (0.4,-0.8325) (0.42,-0.8268) (0.44,-0.8217) (0.46,-0.8169) (0.48,-0.8127) (0.5,-0.809) (0.52,-0.8059) (0.54,-0.8034) (0.56,-0.8015) (0.58,-0.8004) (0.6,-0.8) (0.62,-0.8004) (0.64,-0.8017) (0.66,-0.8038) (0.68,-0.8069) (0.7,-0.811) (0.72,-0.8161) (0.74,-0.8222) (0.76,-0.8294) (0.78,-0.8378) (0.8,-0.8472) (0.82,-0.8578) (0.84,-0.8694) (0.86,-0.8822) (0.88,-0.8961) (0.9,-0.911) (0.92,-0.9269) (0.94,-0.9438) (0.96,-0.9617) (0.98,-0.9804) (1,-1)};
\addplot[cmuted,thin] coordinates {(0,-1) (0.02,-0.9759) (0.04,-0.9518) (0.06,-0.9278) (0.08,-0.9039) (0.1,-0.8801) (0.12,-0.8563) (0.14,-0.8327) (0.16,-0.8092) (0.18,-0.7859) (0.2,-0.7627) (0.22,-0.7397) (0.24,-0.7169) (0.26,-0.6945) (0.28,-0.6723) (0.3,-0.6505) (0.32,-0.6291) (0.34,-0.6082) (0.36,-0.588) (0.38,-0.5685) (0.4,-0.55) (0.42,-0.5326) (0.44,-0.5166) (0.46,-0.5023) (0.48,-0.4903) (0.5,-0.4811) (0.52,-0.4752) (0.54,-0.4734) (0.56,-0.4762) (0.58,-0.484) (0.6,-0.4966) (0.62,-0.5137) (0.64,-0.5346) (0.66,-0.5585) (0.68,-0.5845) (0.7,-0.6121) (0.72,-0.6406) (0.74,-0.6696) (0.76,-0.6988) (0.78,-0.7278) (0.8,-0.7564) (0.82,-0.7845) (0.84,-0.8119) (0.86,-0.8385) (0.88,-0.8643) (0.9,-0.8892) (0.92,-0.9132) (0.94,-0.9362) (0.96,-0.9583) (0.98,-0.9796) (1,-1)};
\addplot[cmuted,thin] coordinates {(0,1) (0.02,0.9559) (0.04,0.9118) (0.06,0.8678) (0.08,0.8239) (0.1,0.7801) (0.12,0.7363) (0.14,0.6927) (0.16,0.6492) (0.18,0.6059) (0.2,0.5627) (0.22,0.5197) (0.24,0.4769) (0.26,0.4345) (0.28,0.3923) (0.3,0.3505) (0.32,0.3091) (0.34,0.2682) (0.36,0.228) (0.38,0.1885) (0.4,0.15) (0.42,0.1126) (0.44,0.07657) (0.46,0.04234) (0.48,0.01033) (0.5,-0.01891) (0.52,-0.04477) (0.54,-0.06657) (0.56,-0.08376) (0.58,-0.09601) (0.6,-0.1034) (0.62,-0.1063) (0.64,-0.1054) (0.66,-0.1015) (0.68,-0.09548) (0.7,-0.08792) (0.72,-0.07941) (0.74,-0.0704) (0.76,-0.06125) (0.78,-0.05224) (0.8,-0.0436) (0.82,-0.03551) (0.84,-0.0281) (0.86,-0.02146) (0.88,-0.01568) (0.9,-0.0108) (0.92,-0.006836) (0.94,-0.003794) (0.96,-0.00166) (0.98,-0.0004078) (1,0)};
\addplot[cmuted,thin] coordinates {(0,1) (0.02,0.9701) (0.04,0.9402) (0.06,0.9105) (0.08,0.8809) (0.1,0.8514) (0.12,0.822) (0.14,0.7928) (0.16,0.7638) (0.18,0.7349) (0.2,0.7062) (0.22,0.6777) (0.24,0.6494) (0.26,0.6213) (0.28,0.5935) (0.3,0.5659) (0.32,0.5386) (0.34,0.5115) (0.36,0.4848) (0.38,0.4585) (0.4,0.4325) (0.42,0.4068) (0.44,0.3817) (0.46,0.3569) (0.48,0.3327) (0.5,0.309) (0.52,0.2859) (0.54,0.2634) (0.56,0.2415) (0.58,0.2204) (0.6,0.2) (0.62,0.1804) (0.64,0.1617) (0.66,0.1438) (0.68,0.1269) (0.7,0.111) (0.72,0.09607) (0.74,0.08222) (0.76,0.06944) (0.78,0.05777) (0.8,0.04721) (0.82,0.03777) (0.84,0.02944) (0.86,0.02222) (0.88,0.01607) (0.9,0.01098) (0.92,0.006905) (0.94,0.003814) (0.96,0.001664) (0.98,0.000408) (1,0)};
\addplot[cmuted,thin] coordinates {(0,1) (0.02,0.9842) (0.04,0.9684) (0.06,0.9526) (0.08,0.937) (0.1,0.9213) (0.12,0.9057) (0.14,0.8901) (0.16,0.8746) (0.18,0.8591) (0.2,0.8436) (0.22,0.8282) (0.24,0.8127) (0.26,0.7973) (0.28,0.7818) (0.3,0.7662) (0.32,0.7506) (0.34,0.735) (0.36,0.7192) (0.38,0.7032) (0.4,0.687) (0.42,0.6706) (0.44,0.6538) (0.46,0.6366) (0.48,0.6189) (0.5,0.6005) (0.52,0.5814) (0.54,0.5614) (0.56,0.5404) (0.58,0.5181) (0.6,0.4945) (0.62,0.4694) (0.64,0.4428) (0.66,0.4146) (0.68,0.3849) (0.7,0.3538) (0.72,0.3217) (0.74,0.2887) (0.76,0.2553) (0.78,0.222) (0.8,0.1892) (0.82,0.1574) (0.84,0.1273) (0.86,0.09931) (0.88,0.07399) (0.9,0.05184) (0.92,0.0333) (0.94,0.01869) (0.96,0.008247) (0.98,0.002036) (1,0)};
\addplot[cmuted,thin] coordinates {(0,3) (0.02,2.93) (0.04,2.86) (0.06,2.791) (0.08,2.722) (0.1,2.653) (0.12,2.584) (0.14,2.516) (0.16,2.448) (0.18,2.38) (0.2,2.313) (0.22,2.246) (0.24,2.179) (0.26,2.113) (0.28,2.047) (0.3,1.982) (0.32,1.918) (0.34,1.854) (0.36,1.792) (0.38,1.729) (0.4,1.668) (0.42,1.608) (0.44,1.549) (0.46,1.492) (0.48,1.436) (0.5,1.381) (0.52,1.329) (0.54,1.278) (0.56,1.23) (0.58,1.185) (0.6,1.142) (0.62,1.103) (0.64,1.067) (0.66,1.035) (0.68,1.006) (0.7,0.9817) (0.72,0.9608) (0.74,0.9437) (0.76,0.9303) (0.78,0.9203) (0.8,0.9137) (0.82,0.9103) (0.84,0.91) (0.86,0.9125) (0.88,0.9179) (0.9,0.9258) (0.92,0.9362) (0.94,0.9489) (0.96,0.9639) (0.98,0.981) (1,1)};
\draw[cink,<->,thick] (axis cs:0.670,-1.690) -- (axis cs:0.670,-1.070)
  node[midway,right,font=\small,fill=white] {$g_{\min}=0.62$};
\node[font=\small,anchor=west] at (axis cs:0.02,-2.75) {$H_B=-\sum_k X_k$ : fondamental $|{+}{+}{+}\rangle$};
\node[font=\small,anchor=east] at (axis cs:0.98,-2.25) {$H_C$ : fondamental $|101\rangle$, $E=-2$};
\end{axis}
\end{tikzpicture}

\caption{Les 8 niveaux d'énergie de $H(s)$ pour le QUBO jouet à 3 qubits (diagonalisation exacte). Le fondamental (bleu) relie $\ket{+++}$ à $\ket{101}$. L'écart minimal avec le premier excité (orange) fixe la durée nécessaire du recuit.}\label{fig:spectre}
\end{figure}

\paragraph{Pourquoi $H_B$ doit-il ne pas commuter avec $H_C$ ?} Si $[H_B,H_C]=0$, les deux auraient les mêmes vecteurs propres, et le système resterait dans $\ket+^{\otimes n}$ sans jamais « choisir » de chaîne. Le terme $X_k$, appelé \emph{champ transverse}, permet aux amplitudes de passer d'une chaîne à ses voisines. C'est l'effet tunnel de la figure~\ref{fig:paysage}.

\paragraph{La limite fondamentale.} Pour les problèmes difficiles, $g_{\min}$ peut décroître \emph{exponentiellement} avec $n$, ce qui impose un temps exponentiel : l'adiabatique ne fait pas de miracle sur les problèmes NP-difficiles.

\section{Le recuit quantique en pratique (D-Wave)}
Les machines D-Wave réalisent physiquement ce programme avec des milliers de qubits supraconducteurs, mais :
\begin{itemize}
\item le système n'est pas isolé (température non nulle, bruit). Ce n'est pas un calcul adiabatique idéal, et le résultat est un \emph{échantillon} de chaînes de basse énergie. On lance de nombreuses lectures (\emph{reads}) ;
\item les qubits ne sont couplés qu'à quelques voisins sur une grille fixe. Un graphe complet comme le nôtre (36 arêtes) exige un \textbf{plongement} (\emph{embedding}) : chaque qubit logique est représenté par une \emph{chaîne} de qubits physiques, maintenus égaux par de forts couplages ferromagnétiques. Les chaînes peuvent se « casser » ;
\item les coefficients $h_i$, $J_{ij}$ ont une plage et une précision limitées. Un QUBO dont les écarts d'énergie utiles sont minuscules devant le plus grand coefficient (pénalités trop fortes, quasi-dégénérescences comme la nôtre) est difficile à résoudre.
\end{itemize}
Le SDK \emph{Ocean} de D-Wave fournit aussi un recuit simulé classique (\texttt{SimulatedAnnealingSampler}), avec la même interface que le matériel : c'est ce que la note d'intention prévoit d'utiliser.

\exercices
\begin{exercice}
Au cours d'un recuit simulé, on propose un mouvement qui augmente l'énergie de $\Delta E=1$. Quelle est la probabilité de l'accepter à $T=2$ ? À $T=0{,}5$ ? À $T=0{,}1$ ? Expliquer le rôle de la température.
\end{exercice}
\begin{exercice}
Pour une famille de problèmes, on mesure $g_{\min}\approx2^{-n/2}$. Comment croît le temps de recuit adiabatique nécessaire avec $n$ ? Comparer au coût de l'énumération exhaustive ($2^n$ évaluations). Qu'en conclure ?
\end{exercice}
